We test the following examples using both directions of each pair of
primaries at equal mixing probability, $p\equiv y=1/2$. For a short
interval $A$ of $\ell$ sites on a periodic chain of $N$ sites, and its
large complement $B$, we write
\begin{align}
 D^A_{a\mid b}(x)&=S\!\left(\rho_a^A\Big\|
                    \frac{\rho_a^A+\rho_b^A}{2}\right),
       &x&=\ell/N,\nonumber\\
 \delta_{a\mid b}(\ep)&=\log2-D^B_{a\mid b},
       &\ep&=1-x_B=\ell/N.
 \label{eq:num-observables}
\end{align}
The order $a\mid b$ specifies the first argument of the relative
entropy; the mixture is the same in both directions. All logarithms
are natural. The large-interval figures retain the label $y=\ell/N$
from the numerical study; this is $\ep$ in the present manuscript,
not the mixing probability.

The calculations use $\ell=6,\ldots,10$, with chains through $N=8192$
for Ising and $N=16384$ for the compact boson. We compare a free
power law, a one-term fit $F_1$ with the theoretical leading power,
and a two-term fit $F_2$ with fixed leading and next powers. The CFT
curves have no fitted parameters. Small-interval fits use $x\leq0.02$
for Ising and $x\leq0.1$ for the boson; deficit fits use
$\ep\leq0.002$ and $\ep\leq0.005$, respectively. Each window is
shared by all six directions and all three fits within its group.
The leading exponents are denoted $\alpha$ for $D^A$ and $\beta$
for $\delta$. All reported amplitudes multiply $x^r$ or $\ep^r$
itself and include the associated powers of $\pi$.
Appendix~\ref{app:num-methods} gives the state constructions,
entropy algorithms, sampling grid, fitting procedure, and numerical
precision checks.
